\documentclass[a4paper,11pt]{ltjsarticle}
\usepackage[haranoaji]{luatexja-preset}
\usepackage{amsmath,amssymb}
\usepackage[top=12mm,bottom=10mm,left=14mm,right=14mm]{geometry}
\usepackage{array}
\usepackage{tikz}
\usetikzlibrary{calc}
\pagestyle{empty}
\setlength{\parindent}{0pt}
\newlength{\qcol}\setlength{\qcol}{0.47\textwidth}
\newlength{\qgap}
\newlength{\anscolw}\setlength{\anscolw}{0.30\textwidth}
\newlength{\ansh}
\newcommand{\Q}[2]{\parbox[t]{\qcol}{\raggedright (#1)\hspace{0.6em}$\displaystyle #2$}}
\newcommand{\QR}[4]{\Q{#1}{#2}\hfill\Q{#3}{#4}\par\vspace{\qgap}}
\newcommand{\anscell}[1]{\rule[-\ansdrop]{0pt}{\ansh}(#1)}
\newlength{\ansdrop}

\begin{document}
{\large\bfseries 数学テスト　9}\hfill 名前(\underline{\hspace{32mm}})\par\vspace{3mm}
■（1）〜（12）の計算をしなさい。（13）、（14）は連立方程式を解きなさい。\par\vspace{3mm}
\setlength{\qgap}{11mm}
\QR{1}{\dfrac{a-b}{4}-\dfrac{a+b}{5}}{2}{9a-\{3b+(7a+b)-5\}}
\QR{3}{(36x-24y)\div6}{4}{9ab^2\div3b^2\times2a}
\QR{5}{(7x-4y)+(-5x+6y)}{6}{(a^2+4a-5)-(-a^2+2a-1)}
\QR{7}{3a\times(-4ab)}{8}{5(2x-y)+3(x-3y)}
\QR{9}{21a^2b\div9ab}{10}{(5x+2y)+(3x-4y)-(x+6y)}
\QR{11}{(-2a)^3}{12}{\dfrac{3}{2}(12x+16y)}
\QR{13}{\left\{\begin{array}{l} 3x-5y=11 \\ x=-2y \end{array}\right.}{14}{\left\{\begin{array}{l} 5x-2y=19 \\ 4x+5y=2 \end{array}\right.}
\vfill
\setlength{\ansh}{12mm}\setlength{\ansdrop}{8mm}%
\begin{center}\begin{tabular}{|p{\anscolw}|p{\anscolw}|p{\anscolw}|}\hline
\anscell{1} & \anscell{2} & \anscell{3} \\\hline
\anscell{4} & \anscell{5} & \anscell{6} \\\hline
\anscell{7} & \anscell{8} & \anscell{9} \\\hline
\anscell{10} & \anscell{11} & \anscell{12} \\\hline
\anscell{13} & \anscell{14} &  \\\hline
\end{tabular}\end{center}

\end{document}
